\documentclass[11pt,letterpaper]{article}
\usepackage[T1]{fontenc}
\usepackage[utf8]{inputenc}
\usepackage{amsmath,amsthm}
\usepackage{newpxtext,newpxmath}
\usepackage[margin=1in,headheight=15pt]{geometry}
\usepackage{microtype,mathtools,booktabs,array}
\usepackage[dvipsnames]{xcolor}
\usepackage{enumitem,fancyhdr,titlesec}
\usepackage[colorlinks=true,linkcolor=MidnightBlue,citecolor=MidnightBlue,urlcolor=MidnightBlue]{hyperref}
\usepackage[nameinlink,noabbrev]{cleveref}
\definecolor{Ink}{HTML}{18374A}
\titleformat{\section}{\large\bfseries\color{Ink}}{\thesection}{0.7em}{}
\titleformat{\subsection}{\normalsize\bfseries\color{Ink}}{\thesubsection}{0.7em}{}
\pagestyle{fancy}
\fancyhf{}
\fancyhead[L]{\small\scshape Ap\'ery arrays and conjectural zeta accelerations}
\fancyhead[R]{\small September 2026}
\fancyfoot[C]{\small\thepage}
\renewcommand{\headrulewidth}{0.3pt}
\setlength{\parskip}{0.3em}
\setlength{\emergencystretch}{2em}
\setcounter{tocdepth}{2}
\numberwithin{equation}{section}
\newtheorem{theorem}{Theorem}[section]
\newtheorem{lemma}[theorem]{Lemma}
\newtheorem{proposition}[theorem]{Proposition}
\newtheorem{corollary}[theorem]{Corollary}
\theoremstyle{definition}
\newtheorem{definition}[theorem]{Definition}
\newtheorem{remark}[theorem]{Remark}
\newcommand{\N}{\mathbb N}
\newcommand{\Z}{\mathbb Z}
\newcommand{\Q}{\mathbb Q}
\newcommand{\Hh}[2]{H_{#1}^{(#2)}}
\newcommand{\ee}{\mathrm e}
\newcommand{\dd}{\mathrm d}
\newcommand{\OEIS}[1]{\href{https://oeis.org/#1}{\texttt{#1}}}
\newcommand{\Tarr}{T}
\newcommand{\Carr}{C}
\newcommand{\zthree}{\zeta(3)}
\newcommand{\ztwo}{\zeta(2)}
\DeclareMathOperator{\sgn}{sgn}
\hypersetup{pdftitle={Apéry Arrays and Conjectural Zeta Accelerations},pdfauthor={Research note prepared with ChatGPT},pdfsubject={Proofs of OEIS A143007 and A108625 conjectural shifted-diagonal formulas}}

\begin{document}
\begin{center}
{\small\scshape A self-contained research note}\par\vspace{0.65em}
{\LARGE\bfseries\color{Ink} Ap\'ery Arrays and Conjectural\par Zeta Accelerations}\par\vspace{0.7em}
{\large Proofs for every shifted diagonal in OEIS A143007,\par and both neighboring diagonals in A108625}\par\vspace{0.8em}
{September 20, 2026}
\end{center}
\vspace{0.5em}
\begin{abstract}
The OEIS array A143007 contains a conjectural one-parameter family of
rapidly convergent series for $\zeta(3)$. Its main diagonal consists of
the classical Ap\'ery numbers. We prove the entire family, not merely a
finite collection of numerical instances. The proof constructs a rational
potential on the nonnegative integer lattice: two explicitly certified
contiguous identities make its edge increments independent of the chosen
path, and elementary estimates identify its limit as $\zeta(3)$.
The same method proves both conjectural neighboring-diagonal formulas
for $\zeta(2)$ in A108625 and extends them to every upper and lower offset.
We also obtain exact finite identities, signed rational error bounds,
all monotone-path expansions, the classical Ap\'ery recurrences, and
sharp fixed-offset asymptotics. The accompanying code checks 22,962
identities in exact arithmetic and independently verifies the finite
symbolic certificates. The analytic proof does not depend on those tests.
\end{abstract}

\noindent\textbf{Status of the claim.}
The targeted formulas are still described conjecturally in the OEIS
versions consulted on the date above. This note supplies a complete proof
of those statements. It does not claim that their underlying Ap\'ery
acceleration mechanism, or every consequence developed here, is new.

\tableofcontents
\newpage

\section{The problem and its mathematical setting}\label{sec:context}

Ap\'ery's sequences are an especially appropriate setting for an OEIS
conjecture: they connect explicit integer sums, second-order recurrences,
rational approximation, and special values of the Riemann zeta function.
The familiar sequence associated with $\zeta(3)$ is
\begin{equation}\label{eq:apery3def}
 A_n=\sum_{j=0}^{n}\binom nj^2\binom{n+j}{j}^2
 =1,5,73,1445,33001,819005,\ldots,
\end{equation}
where the displayed list begins at $n=0$. This is \OEIS{A005259}.
Its $\zeta(2)$ companion, \OEIS{A005258}, is
\begin{equation}\label{eq:apery2def}
 b_n=\sum_{j=0}^{n}\binom nj^2\binom{n+j}{j}
 =1,3,19,147,1251,11253,\ldots.
\end{equation}
These are not incidental integer sequences: they are the denominators
in the classical Ap\'ery approximation constructions
\cite{oeis259,oeis258,cohen,straub}.

The two-variable arrays studied below have these sequences as their
main diagonals. They also have an independent geometric interpretation.
Let
\[
 \mathsf A_m=\{(x_1,\ldots,x_{m+1})\in\Z^{m+1}:\textstyle\sum_i x_i=0\},
 \qquad \|x\|=\frac12\sum_i|x_i|.
\]
The array A108625 counts balls in $\mathsf A_m$; A143007 counts balls
in $\mathsf A_m\times\mathsf A_m$ for the sum of the two norms
\cite{oeis108625,oeis143007,bacher}. Our proof will use finite binomial
sums rather than assume any geometric facts about these lattices.

\subsection{Exactly which conjectures are proved?}

The formula section of A143007 explicitly labels its shifted-diagonal
$\zeta(3)$ formula a ``Conjectural result for other diagonals''
\cite{oeis143007}. We prove it for every integer offset $k\geq0$ in
\cref{thm:zeta3}. The entry attributes its construction to Peter Bala
and dates the original sequence to July 22, 2008.

A108625 describes two neighboring-diagonal $\zeta(2)$ formulas
conjecturally in its comments of July 18, 2008 \cite{oeis108625}.
They are the $k=1$ cases of \cref{thm:zeta2upper,thm:zeta2lower}.
The theorems here also give every other nonnegative offset, with the
necessary distinction between the upper and lower diagonals.

For reproducibility, the consulted internal revisions are A143007,
revision 75, and A108625, revision 98; both carry a May 30, 2026 revision
timestamp. A conjectural label in a database is evidence about the
entry's wording, not a guarantee that no proof exists elsewhere.
Two-dimensional and staircase versions of Ap\'ery acceleration are
classical; Cohen's account explicitly develops that viewpoint
\cite[\S\S3--4]{cohen}. The contribution of this note is a direct,
self-contained proof of the exact OEIS statements, together with explicit
finite certificates and error analysis.

\section{The arrays, their indexing, and the main results}\label{sec:results}

All indices in this note are \emph{square-array} indices. They are not
indices in the one-dimensional antidiagonal flattening used on the
OEIS pages. Set
\begin{equation}\label{eq:L}
 L(m,j)=\binom mj\binom{m+j}{j},\qquad
 L(m,j)=0\quad\hbox{if }j>m,
\end{equation}
and, for $m,n\geq0$, define
\begin{align}
 T_{m,n}&=\sum_{j=0}^{\min(m,n)}L(m,j)L(n,j),\label{eq:T}\\
 C_{m,n}&=\sum_{j=0}^{\min(m,n)}L(m,j)\binom nj.\label{eq:C}
\end{align}
The array $T$ is A143007, and $C$ is A108625, with the dimension parameter
in the first index and the radius in the second. Indeed,
\[
 \binom{m+j}{2j}\binom{2j}{j}=L(m,j),
\]
so \eqref{eq:T} immediately agrees with the OEIS formula. The equivalence
of \eqref{eq:C} with the entry's binomial formula and row generating
function is proved in \cref{app:identification}.

The initial portions are
\[
\begin{array}{c|rrrrr}
T_{m,n}&n=0&1&2&3&4\\\hline
m=0&1&1&1&1&1\\
1&1&5&13&25&41\\
2&1&13&73&253&661\\
3&1&25&253&1445&5741
\end{array}
\qquad
\begin{array}{c|rrrrr}
C_{m,n}&n=0&1&2&3&4\\\hline
m=0&1&1&1&1&1\\
1&1&3&5&7&9\\
2&1&7&19&37&61\\
3&1&13&55&147&309
\end{array}
\]
Thus $T_{m,n}=T_{n,m}$, but generally $C_{m,n}\ne C_{n,m}$.
In particular, transposing $C$ silently would interchange the two
$\zeta(2)$ conjectures and alter their constants and signs.

Both arrays are positive integers and are nondecreasing in each index,
as follows directly from their nonnegative summands. Useful boundary
values are
\begin{align}
 T_{0,n}=T_{m,0}=C_{0,n}=C_{m,0}&=1,\label{eq:boundary}\\
 T_{m,1}=1+2m(m+1),\quad
 C_{m,1}&=1+m(m+1),\quad C_{1,n}=1+2n.\label{eq:firstrow}
\end{align}
Their diagonals are precisely $T_{n,n}=A_n$ and $C_{n,n}=b_n$.

Write
\[
 \Hh{k}{s}=\sum_{j=1}^{k}\frac1{j^s},\qquad
 E_k=2\sum_{j=1}^{k}\frac{(-1)^{j+1}}{j^2},
 \qquad \Hh{0}{s}=E_0=0.
\]

\begin{theorem}[All shifted $\zeta(3)$ diagonals]\label{thm:zeta3}
For every integer $k\geq0$,
\begin{equation}\label{eq:main3}
 \boxed{\displaystyle
 \zthree=\Hh{k}{3}+
 \sum_{n=1}^{\infty}
 \frac{(2n+k)(3n^2+3nk+k^2)}
 {n^3(n+k)^3\,T_{n-1,n+k-1}T_{n,n+k}}.}
\end{equation}
The series is positive and converges geometrically. Its partial sum
through $n=N$ is strictly below $\zthree$.
\end{theorem}

At $k=0$, \eqref{eq:main3} is the classical Ap\'ery series
\begin{equation}\label{eq:classical3}
 \zthree=6\sum_{n=1}^{\infty}\frac1{n^3 A_{n-1}A_n}.
\end{equation}
For example, the first genuinely shifted case is
\[
 \zthree=1+\frac{21}{8\cdot13}
              +\frac{95}{216\cdot13\cdot253}+\cdots.
\]

\begin{theorem}[All upper $\zeta(2)$ diagonals]\label{thm:zeta2upper}
For every integer $k\geq0$,
\begin{equation}\label{eq:main2upper}
 \boxed{\displaystyle
 \ztwo=\Hh{k}{2}+
 \sum_{n=1}^{\infty}(-1)^{n+1}
 \frac{(n+k)^2+(2n+k)^2}
 {n^2(n+k)^2\,C_{n-1,n+k-1}C_{n,n+k}}.}
\end{equation}
The series is absolutely convergent. If $S^{(2,+)}_{N,k}$ denotes its
partial sum including $\Hh{k}{2}$, then
$\sgn(\ztwo-S^{(2,+)}_{N,k})=(-1)^N$.
\end{theorem}

\begin{theorem}[All lower $\zeta(2)$ diagonals]\label{thm:zeta2lower}
For every integer $k\geq0$,
\begin{equation}\label{eq:main2lower}
 \boxed{\displaystyle
 \ztwo=E_k+
 \sum_{n=1}^{\infty}(-1)^{n+k+1}
 \frac{n^2+(2n+k)^2}
 {n^2(n+k)^2\,C_{n+k-1,n-1}C_{n+k,n}}.}
\end{equation}
The series is absolutely convergent. If $S^{(2,-)}_{N,k}$ denotes its
partial sum including $E_k$, then
$\sgn(\ztwo-S^{(2,-)}_{N,k})=(-1)^{N+k}$.
\end{theorem}

At $k=0$ both theorems give
\begin{equation}\label{eq:classical2}
 \ztwo=5\sum_{n=1}^{\infty}\frac{(-1)^{n+1}}{n^2 b_{n-1}b_n}.
\end{equation}
Their $k=1$ cases are exactly the conjectures recorded in A108625:
\begin{align}
 \ztwo&=1+\sum_{n=1}^{\infty}(-1)^{n+1}
 \frac{(n+1)^2+(2n+1)^2}
 {n^2(n+1)^2 C_{n-1,n}C_{n,n+1}},\label{eq:oeis2upper}\\
 \ztwo&=2-\sum_{n=1}^{\infty}(-1)^{n+1}
 \frac{n^2+(2n+1)^2}
 {n^2(n+1)^2 C_{n,n-1}C_{n+1,n}}.\label{eq:oeis2lower}
\end{align}
For comparison with the entry, these begin
\[
 1+\frac{13}{4\cdot5}-\frac{34}{36\cdot5\cdot37}+\cdots,
 \qquad
 2-\frac{10}{4\cdot7}+\frac{29}{36\cdot7\cdot55}-\cdots.
\]

\section{Finite contiguous identities: the algebraic certificates}\label{sec:contiguous}

The proof is driven by identities involving the four corners of one
unit square. There is no experimental recurrence guessing in what follows:
each identity is certified by an explicit telescoping summand.

\subsection{The \texorpdfstring{$\zeta(3)$}{zeta(3)} square}

Fix integers $m,n\geq1$ and abbreviate
\[
 a=T_{m-1,n-1},\quad b=T_{m,n-1},\quad
 c=T_{m-1,n},\quad d=T_{m,n},
 \qquad s=(m+n)(m^2+mn+n^2).
\]

\begin{lemma}\label{lem:Tcell}
The four corners satisfy
\begin{align}
 m^3d+n^3a&=s c,\label{eq:Tcell1}\\
 n^3d+m^3a&=s b.\label{eq:Tcell2}
\end{align}
\end{lemma}
\begin{proof}
Let $M=\min(m,n)$ and $f_j=L(m,j)L(n,j)$ for $0\leq j\leq M$.
The ratios needed to lower the parameters are
\[
 \frac{L(m-1,j)}{L(m,j)}=\frac{m-j}{m+j},\qquad
 \frac{L(n-1,j)}{L(n,j)}=\frac{n-j}{n+j}.
\]
They remain valid at an endpoint where a numerator vanishes.
Define
\[
 R_j=\frac{j^4}{(m+j)(n+j)}f_j\quad(0\leq j\leq M),
 \qquad R_{M+1}=0.
\]
The ratio of consecutive summands gives, for $0\leq j\leq M$,
\begin{equation}\label{eq:Ttelescoper}
 R_{j+1}=(m-j)(n-j)f_j.
\end{equation}
For $j=M$ both sides vanish; for $j<M$ this follows from
\[
 \frac{f_{j+1}}{f_j}
 =\frac{(m-j)(m+j+1)(n-j)(n+j+1)}{(j+1)^4}.
\]
The following rational identity is the certificate:
\begin{align}
 &m^3+n^3\frac{(m-j)(n-j)}{(m+j)(n+j)}
       -s\frac{m-j}{m+j}\notag\\
 &\hspace{2em}=2(m+n)\left\{
 \frac{j^4}{(m+j)(n+j)}-(m-j)(n-j)\right\}.\label{eq:Tcertificate}
\end{align}
It is checked simply by multiplying by $(m+j)(n+j)$ and expanding.
Multiplication by $f_j$ turns its right side into
$2(m+n)(R_j-R_{j+1})$. Summing from $0$ to $M$ proves
\eqref{eq:Tcell1}, since $R_0=R_{M+1}=0$. Interchanging $m$ and $n$
and using the symmetry of $T$ proves \eqref{eq:Tcell2}.
\end{proof}

\subsection{The \texorpdfstring{$\zeta(2)$}{zeta(2)} square}

Now use the same local letters for the corners of $C$:
\[
 a=C_{m-1,n-1},\quad b=C_{m,n-1},\quad
 c=C_{m-1,n},\quad d=C_{m,n},\qquad q=n^2+(m+n)^2.
\]

\begin{lemma}\label{lem:Ccell}
The four corners satisfy
\begin{align}
 2n^2d-m^2a&=q b,\label{eq:Ccell1}\\
 m^2d+2n^2a&=q c.\label{eq:Ccell2}
\end{align}
\end{lemma}
\begin{proof}
Put $g_j=L(m,j)\binom nj$, $M=\min(m,n)$, and set
\[
 r_m=\frac{m-j}{m+j},\qquad r_n=\frac{n-j}{n},\qquad
 R_j=\frac{j^3}{m+j}g_j\quad(0\leq j\leq M),\quad R_{M+1}=0.
\]
Here $r_m$ and $r_n$ are the ratios obtained by lowering $m$ and $n$,
respectively. As in the preceding proof,
\[
 R_{j+1}=(m-j)(n-j)g_j,
\]
because
\[
 \frac{g_{j+1}}{g_j}
 =\frac{(m-j)(m+j+1)(n-j)}{(j+1)^3}
 \quad(j<M).
\]
Two elementary rational identities are
\begin{align}
 m^2+2n^2r_mr_n-q r_m
 &=2\left\{\frac{j^3}{m+j}-(m-j)(n-j)\right\},\label{eq:Ccert2}\\
 2n^2-m^2r_mr_n-q r_n
 &=\frac{2(m+n)}{n}
 \left\{\frac{j^3}{m+j}-(m-j)(n-j)\right\}.\label{eq:Ccert1}
\end{align}
Multiply each identity by $g_j$, sum, and use $R_0=R_{M+1}=0$.
The resulting identities are \eqref{eq:Ccell2} and \eqref{eq:Ccell1}.
All denominators used here are nonzero because $m,n\geq1$.
\end{proof}

These three certificates are sufficient for the whole infinite family
of conjectures. Their role is structural: a local identity will allow
an infinite series to be evaluated by changing its lattice path.

\section{A rational potential proving the \texorpdfstring{$\zeta(3)$}{zeta(3)} conjecture}\label{sec:zeta3proof}

\subsection{Path independence on the integer lattice}

A horizontal edge into $(m,n)$ is the step $(m-1,n)\to(m,n)$, and a
vertical edge into $(m,n)$ is $(m,n-1)\to(m,n)$. If their weights are
$h(m,n)$ and $v(m,n)$, respectively, the square condition is
\begin{equation}\label{eq:closed}
 h(m,n-1)+v(m,n)=v(m-1,n)+h(m,n).
\end{equation}
It says that the two routes from $(m-1,n-1)$ to $(m,n)$ have the same
weight.

\begin{lemma}[Finite path independence]\label{lem:path}
If \eqref{eq:closed} holds for every $m,n\geq1$, then every monotone
lattice path from $(0,0)$ to $(m,n)$ has the same sum of edge weights.
Consequently there is a potential $F_{m,n}$ with $F_{0,0}=0$ and
\[
 F_{m,n}-F_{m-1,n}=h(m,n),\qquad
 F_{m,n}-F_{m,n-1}=v(m,n).
\]
\end{lemma}
\begin{proof}
Encode a monotone path as a word with $m$ horizontal steps and $n$
vertical steps. Adjacent horizontal--vertical pairs can be interchanged
without changing the sum, by \eqref{eq:closed}. Such interchanges connect
any two words with the same numbers of each step. Define $F_{m,n}$ as
this common sum. Appending a last edge proves the asserted increments.
\end{proof}

For $T$, assign the positive rational edge weights
\begin{equation}\label{eq:edges3}
 h_3(m,n)=\frac1{m^3T_{m-1,n}T_{m,n}},\qquad
 v_3(m,n)=\frac1{n^3T_{m,n-1}T_{m,n}}.
\end{equation}
Each expression is used only when the corresponding edge exists.
With the local corner notation of \cref{lem:Tcell}, condition
\eqref{eq:closed}, after clearing positive denominators, is
\[
 c(n^3d+m^3a)=b(m^3d+n^3a).
\]
Both sides are $sbc$, by \cref{lem:Tcell}. Hence the weights are
path independent.

Let their potential be $U_{m,n}$. Taking the two boundary routes yields
the explicit finite identity
\begin{equation}\label{eq:U}
 \boxed{\displaystyle
 U_{m,n}=\Hh{n}{3}+\sum_{i=1}^{m}
   \frac1{i^3T_{i-1,n}T_{i,n}}
 =\Hh{m}{3}+\sum_{j=1}^{n}
   \frac1{j^3T_{m,j-1}T_{m,j}}.}
\end{equation}
In particular $U_{m,n}$ is rational, symmetric, and strictly increasing
when either coordinate increases by one.

\subsection{Identifying the limit without assuming an Ap\'ery theorem}

\begin{lemma}\label{lem:Ulimit}
For all $m,n\geq0$, $U_{m,n}<\zthree$. Moreover,
$U_{m,n}\to\zthree$ whenever $\max(m,n)\to\infty$.
\end{lemma}
\begin{proof}
The second expression in \eqref{eq:U} and monotonicity of $T$ imply
\[
 0\leq U_{m,n}-\Hh{m}{3}
 \leq\frac{\zthree}{T_{m,1}}
 =\frac{\zthree}{1+2m(m+1)}.
\]
This tends to zero as $m\to\infty$, uniformly in $n$; thus for each
fixed $n$, the strictly increasing sequence $U_{m,n}$ tends to $\zthree$.
In particular every finite value is strictly below $\zthree$.
The two expressions in \eqref{eq:U} give
\[
 \Hh{\max(m,n)}{3}\leq U_{m,n}<\zthree.
\]
The last assertion now follows by squeezing.
\end{proof}

It follows, incidentally, that \emph{every} infinite monotone path
starting at $(0,0)$ yields a positive series with sum $\zthree$.
Rows and diagonals are particular paths, not separate approximation
phenomena.

\subsection{The diagonal increment}

The decisive cancellation is finite. A horizontal step followed by a
vertical step, together with \eqref{eq:Tcell2}, gives
\begin{align}
 U_{m,n}-U_{m-1,n-1}
 &=\frac1{m^3ab}+\frac1{n^3bd}\notag\\
 &=\frac{n^3d+m^3a}{m^3n^3abd}\notag\\
 &=\frac{(m+n)(m^2+mn+n^2)}
 {m^3n^3T_{m-1,n-1}T_{m,n}}.\label{eq:diag3}
\end{align}
The intermediate corner $b$ cancels because the contiguous identity
makes its numerator a multiple of $b$. This explains the apparently
special cubic polynomial in the conjectured summand.

\begin{proof}[Proof of \cref{thm:zeta3}]
Apply \eqref{eq:diag3} at the points $(n,n+k)$ and sum the
resulting identity for $n=1,\ldots,N$.
Since $U_{0,k}=\Hh{k}{3}$, the exact result is
\begin{equation}\label{eq:finite3}
 U_{N,N+k}=\Hh{k}{3}+
 \sum_{n=1}^{N}
 \frac{(2n+k)(3n^2+3nk+k^2)}
 {n^3(n+k)^3T_{n-1,n+k-1}T_{n,n+k}}.
\end{equation}
Let $N\to\infty$ and apply \cref{lem:Ulimit}. All terms are positive,
and the partial sums are below the limit. Geometric convergence follows
from the explicit bounds in \cref{sec:errors}, independently of the
sharper asymptotics in \cref{sec:asymptotics}.
\end{proof}

\section{A signed rational potential for the \texorpdfstring{$\zeta(2)$}{zeta(2)} conjectures}\label{sec:zeta2proof}

The array $C$ is not symmetric, and the correct weights necessarily
remember the first coordinate:
\begin{equation}\label{eq:edges2}
 h_2(m,n)=\frac{2(-1)^{m+1}}{m^2C_{m-1,n}C_{m,n}},\qquad
 v_2(m,n)=\frac{(-1)^m}{n^2C_{m,n-1}C_{m,n}}.
\end{equation}
The square condition \eqref{eq:closed}, after clearing denominators and
a common sign, becomes
\[
 c(2n^2d-m^2a)=b(m^2d+2n^2a).
\]
By \cref{lem:Ccell}, both sides are $qbc$. Hence these weights define
a rational potential $V_{m,n}$. The boundary values now differ on the
two axes:
\[
 V_{0,n}=\Hh{n}{2},\qquad V_{m,0}=E_m.
\]
The two boundary routes give
\begin{equation}\label{eq:V}
 \boxed{\displaystyle
 V_{m,n}=\Hh{n}{2}+2\sum_{i=1}^{m}
 \frac{(-1)^{i+1}}{i^2C_{i-1,n}C_{i,n}}
 =E_m+(-1)^m\sum_{j=1}^{n}
 \frac1{j^2C_{m,j-1}C_{m,j}}.}
\end{equation}

\begin{lemma}\label{lem:Vlimit}
The potential $V_{m,n}$ tends to $\ztwo$ whenever
$\max(m,n)\to\infty$.
\end{lemma}
\begin{proof}
First, $E_m\to\ztwo$: separating the even terms in the absolutely
convergent series gives
\[
 \sum_{j\geq1}\frac{(-1)^{j+1}}{j^2}
 =\ztwo-2\sum_{j\geq1}\frac1{(2j)^2}=\frac12\ztwo.
\]
The second expression in \eqref{eq:V} gives the uniform estimate
\begin{equation}\label{eq:Vlimitm}
 |V_{m,n}-E_m|\leq\frac{\ztwo}{C_{m,1}}
 =\frac{\ztwo}{1+m(m+1)}.
\end{equation}
On the other hand, the first expression gives
\begin{equation}\label{eq:Vlimitn}
 |V_{m,n}-\Hh{n}{2}|
 \leq\frac{2\ztwo}{C_{1,n}}=\frac{2\ztwo}{2n+1}.
\end{equation}
For \eqref{eq:Vlimitn}, use $C_{i,n}\geq C_{1,n}$ and
$C_{i-1,n}\geq1$ for every $i\geq1$. Either estimate proves the
limit when its corresponding coordinate tends to infinity, uniformly
in the other. Together they prove the assertion for
$\max(m,n)\to\infty$.
\end{proof}

As before, every infinite monotone path produces a convergent expansion
for the same constant, although the edge weights now have signs.
To combine the two edges of a square, use \eqref{eq:Ccell1}:
\begin{align}
 V_{m,n}-V_{m-1,n-1}
 &=\frac{2(-1)^{m+1}}{m^2ab}+\frac{(-1)^m}{n^2bd}\notag\\
 &=(-1)^{m+1}\frac{2n^2d-m^2a}{m^2n^2abd}\notag\\
 &=(-1)^{m+1}
 \frac{n^2+(m+n)^2}{m^2n^2 C_{m-1,n-1}C_{m,n}}.\label{eq:diag2}
\end{align}

\begin{proof}[Proof of \cref{thm:zeta2upper,thm:zeta2lower}]
Sum \eqref{eq:diag2} along the points $(n,n+k)$, starting from $(0,k)$.
This gives the exact finite identity
\begin{equation}\label{eq:finite2upper}
 V_{N,N+k}=\Hh{k}{2}+\sum_{n=1}^{N}(-1)^{n+1}
 \frac{(n+k)^2+(2n+k)^2}
 {n^2(n+k)^2 C_{n-1,n+k-1}C_{n,n+k}}.
\end{equation}
Instead sum along $(n+k,n)$, starting from $(k,0)$. The result is
\begin{equation}\label{eq:finite2lower}
 V_{N+k,N}=E_k+\sum_{n=1}^{N}(-1)^{n+k+1}
 \frac{n^2+(2n+k)^2}
 {n^2(n+k)^2 C_{n+k-1,n-1}C_{n+k,n}}.
\end{equation}
Take $N\to\infty$ and apply \cref{lem:Vlimit}. Absolute convergence
already follows by bounding the positive arrays below by $1$:
for fixed $k$, the absolute summands are $O(n^{-2})$.
The exact signs of the errors follow from \eqref{eq:error2exact} below.
\end{proof}

The additive constant $E_k$ in the lower formula is not interchangeable
with $\Hh{k}{2}$. Likewise its extra sign $(-1)^k$ is not optional.
Both are forced by the boundary values and the orientation of the
potential.

\section{Exact remainders and certified rational intervals}\label{sec:errors}

A proof of convergence is stronger when it also gives computable error
bounds. Here the potentials make the error an explicit tail in either
coordinate direction, so no numerical evaluation of a zeta function is
required for certification.

\subsection{The positive \texorpdfstring{$\zeta(3)$}{zeta(3)} error}

For fixed $n$, extend the horizontal path from $(m,n)$ to infinity.
By \cref{lem:Ulimit},
\begin{equation}\label{eq:error3exact}
 \zthree-U_{m,n}
 =\sum_{i=m+1}^{\infty}\frac1{i^3T_{i-1,n}T_{i,n}}
 =\sum_{j=n+1}^{\infty}\frac1{j^3T_{m,j-1}T_{m,j}}.
\end{equation}
Both representations are positive. For $m\geq1$, monotonicity gives
\begin{align}
 \frac1{(m+1)^3T_{m,n}T_{m+1,n}}
 &\leq\zthree-U_{m,n}\notag\\
 &\leq\frac{\zthree-\Hh{m}{3}}{T_{m,n}T_{m+1,n}}
 <\frac1{2m^2T_{m,n}T_{m+1,n}}.\label{eq:error3bound}
\end{align}
The last inequality uses
$\sum_{i>m}i^{-3}<\int_m^{\infty}t^{-3}\,\dd t=1/(2m^2)$.
There is a second bound with the coordinates interchanged. Thus, for
$m,n\geq1$, the explicitly rational number
\begin{equation}\label{eq:B3}
 B_3(m,n)=\min\left\{
 \frac1{2m^2T_{m,n}T_{m+1,n}},\,
 \frac1{2n^2T_{m,n}T_{m,n+1}}\right\}
\end{equation}
satisfies
\[
 U_{m,n}<\zthree<U_{m,n}+B_3(m,n).
\]
For the conjectural series, substitute $(m,n)=(N,N+k)$.

\subsection{The signed \texorpdfstring{$\zeta(2)$}{zeta(2)} error}

Taking a vertical tail in \eqref{eq:V} gives the exact identity
\begin{equation}\label{eq:error2exact}
 \ztwo-V_{m,n}=(-1)^m
 \sum_{j=n+1}^{\infty}\frac1{j^2C_{m,j-1}C_{m,j}}.
\end{equation}
Every summand on the right after removal of $(-1)^m$ is strictly
positive. Hence the sign is exactly $(-1)^m$, not merely eventually so.
For $n\geq1$, the vertical estimates are
\begin{equation}\label{eq:error2vertical}
 \frac1{(n+1)^2C_{m,n}C_{m,n+1}}
 \leq |\ztwo-V_{m,n}|
 <\frac1{nC_{m,n}C_{m,n+1}}.
\end{equation}
The upper bound follows from $\sum_{j>n}j^{-2}<1/n$.

There is a complementary horizontal tail:
\begin{equation}\label{eq:error2horizontal}
 \ztwo-V_{m,n}
 =2\sum_{i=m+1}^{\infty}
 \frac{(-1)^{i+1}}{i^2C_{i-1,n}C_{i,n}}.
\end{equation}
Its magnitudes strictly decrease to zero, since $i^2$ strictly increases
and both array factors are nondecreasing. The alternating-series bound
therefore gives
\begin{equation}\label{eq:error2hbound}
 |\ztwo-V_{m,n}|
 \leq\frac2{(m+1)^2C_{m,n}C_{m+1,n}}.
\end{equation}
For $n\geq1$, define
\begin{equation}\label{eq:B2}
 B_2(m,n)=\min\left\{
 \frac2{(m+1)^2C_{m,n}C_{m+1,n}},\,
 \frac1{nC_{m,n}C_{m,n+1}}\right\}.
\end{equation}
Then $V_{m,n}$ and $V_{m,n}+(-1)^m B_2(m,n)$ are endpoints of a
closed rational interval containing $\ztwo$. The parity gives the
correct ordering of its endpoints. The same construction works when
$n=0$ by using only the first term in \eqref{eq:B2}.

\subsection{An elementary proof of geometric acceleration}

The bounds above already imply geometric convergence, without any
saddle-point analysis. Since the central binomial coefficient is the
largest of the $2N+1$ coefficients whose sum is $4^N$,
\[
 \binom{2N}{N}\geq\frac{4^N}{2N+1}.
\]
The $j=N$ summand in $T_{N,N+k}$ gives
\[
 T_{N,N+k}\geq\binom{2N}{N}^2\geq\frac{16^N}{(2N+1)^2}.
\]
Apply \eqref{eq:error3bound} and monotonicity to obtain, for $N\geq1$,
\begin{equation}\label{eq:rough3}
 0<\zthree-U_{N,N+k}
 <\frac{(2N+1)^4}{2N^2\,256^N}.
\end{equation}
Likewise, both $C_{N,N+k}$ and $C_{N+k,N}$ are at least
$\binom{2N}{N}$. Using \eqref{eq:error2hbound} gives
\begin{equation}\label{eq:rough2}
 \max\{|\ztwo-V_{N,N+k}|,|\ztwo-V_{N+k,N}|\}
 \leq\frac{2(2N+1)^2}{(N+1)^2\,16^N}.
\end{equation}
These deliberately simple bounds are uniform over all offsets $k\geq0$.
The sharp exponential rates for fixed $k$ are substantially faster.

\subsection{Numerical certificates, not merely decimal agreement}

\cref{tab:bounds} reports rational bounds \eqref{eq:B3} and \eqref{eq:B2}
computed by exact integer arithmetic. The displayed scientific notation
has been rounded \emph{upward}; the exact fractions are included in
\texttt{data/certified\_intervals.csv}. The full CSV also records the
rational partial sums, the error signs, and decimal displays.

\begin{table}[htbp]
\centering
\small
\caption{Rigorous absolute-error upper bounds after $N$ diagonal terms.
Each displayed value is rounded upward to three significant digits.}
\label{tab:bounds}
\begin{tabular}{llrrr}
\toprule
Family & $k$ & $N=10$ & $N=20$ & $N=40$ \\
\midrule
$\zeta(3)$ & 0 & $5.27\cdot10^{-30}$ & $2.26\cdot10^{-60}$ & $2.44\cdot10^{-121}$ \\
 & 1 & $1.79\cdot10^{-31}$ & $7.15\cdot10^{-62}$ & $7.45\cdot10^{-123}$ \\
 & 3 & $3.99\cdot10^{-34}$ & $1.05\cdot10^{-64}$ & $8.48\cdot10^{-126}$ \\
\addlinespace
$\zeta(2)$, upper & 0 & $4.57\cdot10^{-21}$ & $5.88\cdot10^{-42}$ & $9.46\cdot10^{-84}$ \\
 & 1 & $7.18\cdot10^{-22}$ & $8.91\cdot10^{-43}$ & $1.41\cdot10^{-84}$ \\
 & 3 & $2.69\cdot10^{-23}$ & $2.59\cdot10^{-44}$ & $3.53\cdot10^{-86}$ \\
\addlinespace
$\zeta(2)$, lower & 0 & $4.57\cdot10^{-21}$ & $5.88\cdot10^{-42}$ & $9.46\cdot10^{-84}$ \\
 & 1 & $2.80\cdot10^{-22}$ & $3.44\cdot10^{-43}$ & $5.40\cdot10^{-85}$ \\
 & 3 & $1.64\cdot10^{-24}$ & $1.50\cdot10^{-45}$ & $2.00\cdot10^{-87}$ \\
\bottomrule
\end{tabular}

\end{table}

For instance, the unshifted $\zeta(3)$ series after 40 terms has
absolute error less than $2.44\cdot10^{-121}$. This is an absolute-error
statement; claiming that a certain decimal string is correctly rounded
would additionally require testing the certified interval against the
relevant rounding boundary. The archive preserves exact fractions so
that such tests can be made without relying on a floating-point zeta
implementation.

\section{Recovering the classical Ap\'ery recurrences}\label{sec:recurrences}

The square identities also explain why these arrays contain the familiar
Ap\'ery denominators. This provides an independent consistency check and
connects the proof with the usual one-dimensional formulation.

\subsection{Denominator recurrences}

On the diagonal of \cref{lem:Tcell},
\begin{equation}\label{eq:Tneighbor}
 T_{n,n-1}=T_{n-1,n}=\frac{A_n+A_{n-1}}6\qquad(n\geq1).
\end{equation}
At the square with upper-right corner $(n+1,n)$,
\eqref{eq:Tcell1} gives
\[
 (n+1)^3T_{n+1,n}+n^3T_{n,n-1}
 =(2n+1)(3n^2+3n+1)A_n.
\]
Substitute \eqref{eq:Tneighbor}. After multiplication by $6$ and
collection of terms, the result is
\begin{equation}\label{eq:rec3}
 (n+1)^3 A_{n+1}
 =(34n^3+51n^2+27n+5)A_n-n^3A_{n-1},
 \qquad n\geq1.
\end{equation}
The initial values are $A_0=1$ and $A_1=5$.

Similarly, the two diagonal identities for $C$ give
\begin{equation}\label{eq:Cneighbors}
 C_{n,n-1}=\frac{2b_n-b_{n-1}}5,\qquad
 C_{n-1,n}=\frac{b_n+2b_{n-1}}5.
\end{equation}
Use \eqref{eq:Ccell2} at $(m,n)=(n+1,n)$:
\[
 (n+1)^2C_{n+1,n}+2n^2C_{n,n-1}
 =\bigl(n^2+(2n+1)^2\bigr)b_n.
\]
Substitution from \eqref{eq:Cneighbors} gives
\begin{equation}\label{eq:rec2}
 (n+1)^2b_{n+1}=(11n^2+11n+3)b_n+n^2b_{n-1},
 \qquad n\geq1,
\end{equation}
with $b_0=1$ and $b_1=3$. These are the classical recurrences recorded
for A005259 and A005258 \cite{oeis259,oeis258}.

\subsection{The companion numerator solutions}

Define rational arrays
\begin{equation}\label{eq:numerators}
 P_{m,n}=T_{m,n}U_{m,n},\qquad R_{m,n}=C_{m,n}V_{m,n}.
\end{equation}
They satisfy exactly the same local linear identities as $T$ and $C$,
respectively. For example, abbreviate the four $U$ values by
$U_a,U_b,U_c,U_d$. Then
\begin{align*}
 m^3dU_d+n^3aU_a-scU_c
 &=m^3d(U_d-U_c)+n^3a(U_a-U_c)\\
 &=m^3d\,h_3(m,n)-n^3a\,v_3(m-1,n)\\
 &=\frac1c-\frac1c=0.
\end{align*}
The other identity follows by symmetry.
For $R$, the corresponding cancellations are
\begin{align*}
 2n^2dV_d-m^2aV_a-qbV_b
 &=2n^2d\,v_2(m,n)+m^2a\,h_2(m,n-1)=0,\\
 m^2dV_d+2n^2aV_a-qcV_c
 &=m^2d\,h_2(m,n)-2n^2a\,v_2(m-1,n)=0.
\end{align*}
Thus the same elimination used for the denominators applies to their
numerator companions.

Set $p_n=P_{n,n}$ and $r_n=R_{n,n}$. Then $p_n$ satisfies
\eqref{eq:rec3}, with $p_0=0,p_1=6$, and $r_n$ satisfies
\eqref{eq:rec2}, with $r_0=0,r_1=5$. In particular,
\[
 \frac{p_n}{A_n}=U_{n,n}\longrightarrow\zthree,\qquad
 \frac{r_n}{b_n}=V_{n,n}\longrightarrow\ztwo.
\]
The adjacent determinant identities (often called Casoratians) are
immediate from the diagonal increments:
\begin{align}
 p_nA_{n-1}-p_{n-1}A_n&=\frac6{n^3},\label{eq:cas3}\\
 r_nb_{n-1}-r_{n-1}b_n&=\frac{5(-1)^{n+1}}{n^2}.\label{eq:cas2}
\end{align}
The normalization here allows rational numerator sequences; no claim
that all $p_n$ or $r_n$ are integers is intended. In the standard
irrationality arguments, additional denominator-control statements are
important. They are not needed to prove any of the series identities
in this note.

\section{Sharp growth and error asymptotics}\label{sec:asymptotics}

The previous sections already prove the conjectures and give rigorous
rational error bounds. We now determine the exact fixed-offset
asymptotics, both to quantify the acceleration and to distinguish the
effect of an offset from the effect of taking more diagonal terms.

Put
\begin{align}
 \tau&=1+\sqrt2,& q_3&=\tau^2,&
 \lambda_3&=\tau^4=17+12\sqrt2,\label{eq:constants3}\\
 \varphi&=\frac{1+\sqrt5}{2},&&&
 \lambda_2&=\varphi^5=\frac{11+5\sqrt5}{2},\label{eq:constants2}\\
 c_3&=\frac{q_3}{2^{9/4}\pi^{3/2}},&&&
 c_2&=\frac{\varphi^{5/2}}{2\pi\,5^{1/4}}.\label{eq:cconstants}
\end{align}

\begin{theorem}[Fixed-offset array growth]\label{thm:arrayasymp}
For every fixed integer $k\geq0$, as $N\to\infty$,
\begin{align}
 T_{N,N+k}&\sim c_3q_3^k\lambda_3^N N^{-3/2},\label{eq:Tasymp}\\
 C_{N,N+k}&\sim c_2\varphi^{2k}\lambda_2^N N^{-1},\label{eq:Casympupper}\\
 C_{N+k,N}&\sim c_2\varphi^{3k}\lambda_2^N N^{-1}.\label{eq:Casymplower}
\end{align}
\end{theorem}

\subsection{A discrete Laplace estimate}

We use the following elementary form of the saddle-point principle.
Suppose $f$ has a unique interior maximum at $t_0\in(0,1)$,
$f''(t_0)<0$, and the summands near $j/N=t_0$ have the uniform form
\[
 N^{-a}g(j/N)\ee^{Nf(j/N)}(1+o(1)),\qquad g(t_0)>0.
\]
If summands outside every fixed neighborhood of $t_0$ are bounded by a
polynomial in $N$ times $\ee^{N(f(t_0)-\delta)}$ for some $\delta>0$,
then their sum is asymptotic to
\begin{equation}\label{eq:laplace}
 g(t_0)\sqrt{\frac{2\pi}{|f''(t_0)|}}\,
 N^{1/2-a}\ee^{Nf(t_0)}.
\end{equation}
To see this, expand
$f(t)=f(t_0)+\tfrac12f''(t_0)(t-t_0)^2+O(|t-t_0|^3)$ and use
$x=(j-Nt_0)/\sqrt N$. The local sum becomes a Riemann sum for the
Gaussian integral with mesh $N^{-1/2}$; the complement is negligible
by the strict maximum and a Gaussian majorant. First restrict to a
fixed bounded $x$ interval, then let that interval expand. This also
justifies multiplying the summands by a fixed continuous function of
$j/N$: the leading constant is multiplied by its value at $t_0$.

\subsection{The \texorpdfstring{$T$}{T} saddle}

For $j/N=t$ in a compact subinterval of $(0,1)$, Stirling's formula gives
\begin{equation}\label{eq:Lstirling}
 L(N,j)\sim\frac1{2\pi Nt}
 \sqrt{\frac{1+t}{1-t}}\,\ee^{N\ell(t)},
\end{equation}
where
\[
 \ell(t)=(1+t)\log(1+t)-(1-t)\log(1-t)-2t\log t.
\]
The exponent for $T_{N,N}$ is $f_3(t)=2\ell(t)$, and its amplitude
is $N^{-2}g_3(t)$ with
\[
 g_3(t)=\frac{1+t}{4\pi^2t^2(1-t)}.
\]
Since
\[
 \ell'(t)=\log\frac{1-t^2}{t^2},
\]
the unique maximum is at $t_3=1/\sqrt2$. Direct substitution gives
\[
 \ee^{f_3(t_3)}=\lambda_3,\quad
 f_3''(t_3)=-8\sqrt2,\quad
 g_3(t_3)=\frac{q_3}{2\pi^2}.
\]
Equation \eqref{eq:laplace} now gives $c_3\lambda_3^N N^{-3/2}$.

For a fixed offset, the exact ratio is
\begin{equation}\label{eq:Loffset}
 \frac{L(N+k,j)}{L(N,j)}
 =\prod_{r=1}^{k}\frac{N+j+r}{N-j+r}.
\end{equation}
Near $t_3$, this tends uniformly to
$((1+t)/(1-t))^k$, whose value at $t_3$ is $q_3^k$.
This proves \eqref{eq:Tasymp}.

For completeness, the endpoint estimates required by \eqref{eq:laplace}
follow directly from factorial bounds for the binomial summands.
Their logarithms equal $Nf_3(j/N)+O(\log N)$ uniformly when the entropy
terms are extended continuously to $0$ and $1$. For fixed $k$, the
ratio \eqref{eq:Loffset} is bounded by a polynomial in $N$ on the
whole summation range. Since $f_3$ has a strict interior maximum,
the endpoints and every region separated from $t_3$ are exponentially
smaller than the saddle contribution.

\subsection{The \texorpdfstring{$C$}{C} saddle and its asymmetry}

Multiplying \eqref{eq:Lstirling} by Stirling's approximation for
$\binom Nj$ gives exponent
\[
 f_2(t)=(1+t)\log(1+t)-2(1-t)\log(1-t)-3t\log t
\]
and amplitude $N^{-3/2}g_2(t)$, where
\[
 g_2(t)=\frac{\sqrt{1+t}}{(2\pi)^{3/2}t^{3/2}(1-t)}.
\]
The stationary equation is
\[
 f_2'(t)=\log\frac{(1+t)(1-t)^2}{t^3}=0,
 \qquad 1-t-t^2=0.
\]
Thus the unique maximum occurs at $t_2=1/\varphi$. At that point,
\[
 \ee^{f_2(t_2)}=\lambda_2,\qquad
 f_2''(t_2)=-\sqrt5\,\varphi^3,
\]
and \eqref{eq:laplace} gives $c_2\lambda_2^N N^{-1}$.

For the upper diagonal, the offset multiplies the summand by
\[
 \frac{\binom{N+k}{j}}{\binom Nj}
 =\prod_{r=1}^{k}\frac{N+r}{N-j+r}
 \longrightarrow(1-t)^{-k},
\]
which equals $\varphi^{2k}$ at $t_2$. For the lower diagonal the
multiplier is \eqref{eq:Loffset}, with limiting value
$((1+t_2)/(1-t_2))^k=\varphi^{3k}$. The same endpoint argument as
for $T$ applies. This proves \cref{thm:arrayasymp}, including the
asymmetric constants in its last two formulas.

\subsection{Exact leading constants for the errors}

Let $S^{(3)}_{N,k}=U_{N,N+k}$, and retain the notation
$S^{(2,+)}_{N,k}=V_{N,N+k}$ and $S^{(2,-)}_{N,k}=V_{N+k,N}$.

\begin{theorem}[Sharp fixed-offset errors]\label{thm:errorasymp}
For each fixed $k\geq0$,
\begin{align}
 \zthree-S^{(3)}_{N,k}
 &\sim\frac{96\sqrt2\,\pi^3}{\lambda_3^k(\lambda_3^2-1)}
       \lambda_3^{-2N},\label{eq:errorasymp3}\\
 \ztwo-S^{(2,+)}_{N,k}
 &\sim(-1)^N
 \frac{20\pi^2\sqrt5}{\varphi^{4k}(\lambda_2^2+1)}
       \lambda_2^{-2N},\label{eq:errorasymp2upper}\\
 \ztwo-S^{(2,-)}_{N,k}
 &\sim(-1)^{N+k}
 \frac{20\pi^2\sqrt5}{\varphi^{6k}(\lambda_2^2+1)}
       \lambda_2^{-2N}.\label{eq:errorasymp2lower}
\end{align}
In the signed formulas, $\sim$ means that the ratio to the displayed
signed expression tends to $1$.
\end{theorem}
\begin{proof}
Let $a_{n,k}$ be the positive summand in \eqref{eq:main3}.
By \cref{thm:arrayasymp},
\[
 a_{n,k}\sim \frac{6\lambda_3}{c_3^2q_3^{2k}}\lambda_3^{-2n}
 =\frac{96\sqrt2\,\pi^3}{\lambda_3^k}\lambda_3^{-2n}.
\]
The powers $n^{-3}$ in the rational prefactor and in the product of
array asymptotics cancel. Summing the resulting geometric tail gives
\eqref{eq:errorasymp3}.

For the upper $\zeta(2)$ series, the absolute value of the $n$th
summand is asymptotic to
\[
 \frac{5\lambda_2}{c_2^2\varphi^{4k}}\lambda_2^{-2n}
 =\frac{20\pi^2\sqrt5}{\varphi^{4k}}\lambda_2^{-2n}.
\]
For the lower series, replace $\varphi^{4k}$ by $\varphi^{6k}$ and
include the additional sign $(-1)^k$. An alternating geometric tail
has denominator $1+\lambda_2^{-2}$, which yields
\eqref{eq:errorasymp2upper} and \eqref{eq:errorasymp2lower}.

To justify these tail operations directly, if $a_n/(K\rho^n)\to1$
with $0<\rho<1$, then
$\sup_{n>N}|a_n/(K\rho^n)-1|\to0$. The absolute error incurred by
replacing the entire tail with the corresponding geometric tail is
therefore $o(\rho^N)$, in both the positive and alternating cases.
No exchange of an uncontrolled asymptotic expansion with a sum is used.
\end{proof}

The asymptotic decimal gain per extra diagonal term is
\[
 2\log_{10}\lambda_3=3.0622054827\ldots
 \quad\hbox{for }\zthree,
 \qquad
 2\log_{10}\lambda_2=2.0898764025\ldots
 \quad\hbox{for }\ztwo.
\]
A fixed offset changes the leading constant, not these exponential
rates. Increasing $k$ also increases the size and cost of the array
entries, so the prefactor improvement alone does not establish an
optimal computational strategy.

\subsection{Why a fixed row is slower than a diagonal}

For comparison, fix $m$ and let $n\to\infty$. The leading binomial
terms give
\[
 T_{m,n}\sim\alpha_m n^{2m},\quad
 \alpha_m=\frac{\binom{2m}{m}}{(m!)^2};
 \qquad
 C_{m,n}\sim\beta_m n^m,\quad
 \beta_m=\frac{\binom{2m}{m}}{m!}.
\]
Using the exact vertical tails and the elementary asymptotic
$\sum_{j>n}j^{-s}\sim n^{1-s}/(s-1)$ for $s>1$, we obtain
\begin{align}
 \zthree-U_{m,n}&\sim
 \frac{n^{-(4m+2)}}{\alpha_m^2(4m+2)},\label{eq:rowasymp3}\\
 \ztwo-V_{m,n}&\sim(-1)^m
 \frac{n^{-(2m+1)}}{\beta_m^2(2m+1)}.\label{eq:rowasymp2}
\end{align}
Thus a fixed row gives increasing algebraic acceleration as the row
index grows, whereas a fixed-offset diagonal gives geometric
acceleration. Both behaviors are consequences of the same potential.

\section{Reproducible verification and computational use}\label{sec:verification}

The proof is finite algebra followed by explicit convergence estimates.
The accompanying computations are a separate audit of indexing,
normalization, and signs. In particular, the proof of an infinite
identity is not inferred from decimal agreement.

\subsection{What was checked}

The program \texttt{verify.py} uses only Python's standard library, with
arbitrary-precision integers and rational fractions. Its default run
checks the following exact equalities:
\begin{center}
\begin{tabular}{lr}
\toprule
Test group & Equalities\\
\midrule
Square identities, closed edges, and diagonal increments &12,800\\
Potentials, symmetry, and independent array formulas &8,405\\
Finite diagonal formulas, offsets $0\leq k\leq12$ &1,599\\
Denominator and numerator recurrences &156\\
Published diagonal prefixes &2\\
\midrule
Total &22,962\\
\bottomrule
\end{tabular}
\end{center}
The square tests use $1\leq m,n\leq40$. The potential and independent
array tests use $0\leq m,n\leq40$. Finite diagonal formulas are checked
for every $0\leq N\leq40$ and $0\leq k\leq12$, including all three
orientations/families and the empty-sum cases.

For additional independence, $C$ is compared with the coefficient
formula proved in \cref{app:identification}, and $T$ is compared with
the alternative OEIS sum
\[
 \sum_{j=0}^{\min(m,n)}\binom mj^2\binom{m+n-j}{m}^2
\]
\cite{oeis143007}. These comparisons are finite implementation tests;
the proof of the conjectures uses \eqref{eq:T} and \eqref{eq:C} directly.

The optional program \texttt{symbolic\_certificates.py} uses SymPy to
reduce the three rational certificates
\eqref{eq:Tcertificate}, \eqref{eq:Ccert2}, \eqref{eq:Ccert1}, and the
two recurrence-elimination coefficients, to zero. All five tests pass.
The supplied record identifies the tested SymPy version. These are
symbolic identities, not evaluations at sample numerical points.

\subsection{How to reproduce the artifacts}

From the archive directory, run
\begin{verbatim}
python verify.py --max-index 40 --offsets 12 --output-dir data
python make_error_table.py
python symbolic_certificates.py
pdflatex -interaction=nonstopmode -halt-on-error article.tex
pdflatex -interaction=nonstopmode -halt-on-error article.tex
\end{verbatim}
The symbolic command requires SymPy; it is not required for the first
two commands or for compiling the supplied article. The optional
\texttt{numerical\_checks.py} compares the exact partial sums with
high-precision zeta values and illustrates the asymptotic constants.
Its output is not used as an error certificate.

The CSV interval data store numerators and denominators rather than
only decimal approximations. For $\zeta(3)$, use the endpoints
$S$ and $S+B$. For $\zeta(2)$, the recorded sign $\sigma\in\{1,-1\}$
gives endpoints $S$ and $S+\sigma B$, ordered increasingly.
The article, source notes, test programs, and generated data are all
included in the archive.

\subsection{A practical evaluation route}

A direct array entry is a finite sum of $\min(m,n)+1$ positive integers.
The summand ratios in \cref{sec:contiguous} allow successive summands to
be generated without repeatedly evaluating factorials. For a fixed
diagonal and $N$ terms, the direct implementation uses $O(N^2)$
summand-generation steps when $k$ is fixed; this counts arithmetic
operations, not bit complexity. The integers grow exponentially in
value and hence linearly in bit length along a fixed-offset diagonal.

For the main diagonals, \eqref{eq:rec3} and \eqref{eq:rec2} reduce
denominator generation to $O(N)$ recurrence steps. The same recurrences
for $p_n$ and $r_n$ produce the associated rational approximants.
Whichever method is used, one can stop when the rational bound is below
a prescribed absolute tolerance. The positive binomial formulas are
particularly convenient for independent verification and avoid relying
on cancellation in a recurrence.

\section{Conclusion and scope}\label{sec:conclusion}

The conjectural family in A143007 is true for every nonnegative offset.
Both neighboring-diagonal conjectures in A108625 are also true, and
belong to two full offset families. The proof does not assume any of the
conjectural formulas, a known diagonal limit, or Ap\'ery's irrationality
theorem. Instead it identifies the limit from the ordinary convergent
series for $\zeta(3)$ and $\zeta(2)$ on the coordinate axes.

The structural reason is a closed rational system of edge increments.
A finite telescoping certificate yields each square identity; the square
identities yield path independence; a two-edge diagonal step cancels its
intermediate array entry; and the boundary limits identify the sum.
This chain of implications explains the cubic and quadratic numerators
in the OEIS formulas, as well as the differing constants and signs in
the two $\zeta(2)$ orientations.

The conclusions concern the stated acceleration formulas; no new
irrationality theorem or resolution of unrelated Ap\'ery-polynomial
conjectures is claimed.

\clearpage
\appendix
\section{Identification with the OEIS arrays and row generating functions}\label{app:identification}

This appendix verifies the array identification algebraically and
explains the Legendre-polynomial structure behind the crystal-ball
interpretation. It is independent of the potential proof.

\subsection{Two finite forms of the Legendre polynomial}

Define the Legendre polynomial by Rodrigues' formula,
\[
 P_m(z)=\frac1{2^m m!}\frac{\dd^m}{\dd z^m}(z^2-1)^m.
\]
Expansion at $z=1$ gives
\begin{equation}\label{eq:LegendreL}
 P_m(1+2u)=\sum_{j=0}^{m}L(m,j)u^j.
\end{equation}
Indeed, after substituting $z=1+2u$, Rodrigues' formula becomes
$m!^{-1}(\dd/\dd u)^m[u^m(1+u)^m]$, whose coefficient of $u^j$ is
$\binom mj\binom{m+j}{j}$.

Applying the Leibniz rule instead gives
\begin{equation}\label{eq:LegendreW}
 P_m(z)=2^{-m}\sum_{i=0}^{m}\binom mi^2
 (z-1)^i(z+1)^{m-i}.
\end{equation}
With $W_m(x)=\sum_{i=0}^{m}\binom mi^2x^i$, this implies
\begin{equation}\label{eq:Psub}
 P_m\!\left(\frac{1+x}{1-x}\right)=\frac{W_m(x)}{(1-x)^m}.
\end{equation}

\subsection{The row generating function of \texorpdfstring{$C$}{C}}

Using $\sum_{n\geq0}\binom njx^n=x^j/(1-x)^{j+1}$ in
\eqref{eq:C}, we obtain, as a formal power series,
\begin{align}
 \sum_{n\geq0}C_{m,n}x^n
 &=\frac1{1-x}\sum_{j=0}^{m}L(m,j)
       \left(\frac{x}{1-x}\right)^j\notag\\
 &=\frac1{1-x}P_m\!\left(\frac{1+x}{1-x}\right)
 =\frac{W_m(x)}{(1-x)^{m+1}}.\label{eq:Cgf}
\end{align}
Extracting the coefficient of $x^n$ gives
\begin{equation}\label{eq:Calternative}
 C_{m,n}=\sum_{i=0}^{\min(m,n)}\binom mi^2\binom{m+n-i}{m}.
\end{equation}
This is exactly the defining finite sum and row generating function
listed in A108625 \cite{oeis108625}. In particular, the orientation of
the two indices has been verified, rather than inferred from a short
sequence prefix.

\subsection{A finite square identity and the row generating function of \texorpdfstring{$T$}{T}}

We next prove the finite identity
\begin{equation}\label{eq:clausen}
 P_m(z)^2=\sum_{j=0}^{m}L(m,j)\binom{2j}{j}
       \left(\frac{z^2-1}{4}\right)^j.
\end{equation}
It is a terminating Legendre square identity. A short differential
verification is included to keep the identification self-contained.

Put $M=m(m+1)$ and $Q(z)=P_m(z)^2$. The Legendre equation
$(1-z^2)P_m''-2zP_m'+MP_m=0$, obtained by differentiating Rodrigues'
formula, yields the following equation by the product rule:
\begin{align}
 &(z^2-1)^2Q'''+6z(z^2-1)Q''\notag\\
 &\qquad+\bigl((6-4M)z^2+4M-2\bigr)Q'-4MzQ=0.\label{eq:symmetric_square}
\end{align}
Since $P_m(-z)=(-1)^mP_m(z)$, $Q$ is a polynomial in $z^2$.
Set $t=(z^2-1)/4$ and write $Q(z)=y(t)$. Substitution in
\eqref{eq:symmetric_square} and cancellation as a polynomial identity
gives
\begin{equation}\label{eq:yode}
 t^2(4t+1)y'''+(18t^2+3t)y''+
 \bigl((12-4M)t+1\bigr)y'-2My=0.
\end{equation}
For $y(t)=\sum_{j\geq0}q_jt^j$, coefficient comparison gives
\[
 (j+1)^3q_{j+1}=2(2j+1)(m-j)(m+j+1)q_j,
 \qquad q_0=P_m(1)^2=1.
\]
The coefficients $q_j=L(m,j)\binom{2j}{j}$ satisfy exactly this
recurrence and initial value, and vanish for $j>m$. This proves
\eqref{eq:clausen}.

Now
\[
 \sum_{n\geq0}L(n,j)x^n
 =\binom{2j}{j}\frac{x^j}{(1-x)^{2j+1}},
\]
so \eqref{eq:T} and \eqref{eq:clausen} imply
\begin{align}
 \sum_{n\geq0}T_{m,n}x^n
 &=\frac1{1-x}\sum_{j=0}^{m}L(m,j)\binom{2j}{j}
        \left(\frac{x}{(1-x)^2}\right)^j\notag\\
 &=\frac1{1-x}P_m\!\left(\frac{1+x}{1-x}\right)^2
 =\frac{W_m(x)^2}{(1-x)^{2m+1}}.\label{eq:Tgf}
\end{align}
These are the row generating functions recorded in A143007.
The passage from $W_m(x)/(1-x)^m$ to its square is exactly what the
sum norm on a product lattice requires for sphere counts; division by
$1-x$ then converts sphere counts to ball counts
\cite{oeis143007,bacher}.

\clearpage
\begin{thebibliography}{9}
\addcontentsline{toc}{section}{References}

\bibitem{oeis143007}
OEIS Foundation Inc., \emph{The On-Line Encyclopedia of Integer
Sequences}, A143007, square array for the crystal-ball sequences of
$\mathsf A_n\times\mathsf A_n$, originally contributed by Peter Bala,
July 22, 2008. Internal revision 75, May 30, 2026. Consulted September
20, 2026. \url{https://oeis.org/A143007/internal}.

\bibitem{oeis108625}
OEIS Foundation Inc., \emph{The On-Line Encyclopedia of Integer
Sequences}, A108625, square array for the crystal-ball sequences of
$\mathsf A_n$; relevant zeta-acceleration comments by Peter Bala,
July 18, 2008. Internal revision 98, May 30, 2026. Consulted September
20, 2026. \url{https://oeis.org/A108625/internal}.

\bibitem{oeis259}
OEIS Foundation Inc., \emph{The On-Line Encyclopedia of Integer
Sequences}, A005259, Ap\'ery numbers
$\sum_{j=0}^n\binom nj^2\binom{n+j}{j}^2$.
Consulted September 20, 2026. \url{https://oeis.org/A005259}.

\bibitem{oeis258}
OEIS Foundation Inc., \emph{The On-Line Encyclopedia of Integer
Sequences}, A005258, Ap\'ery numbers
$\sum_{j=0}^n\binom nj^2\binom{n+j}{j}$.
Consulted September 20, 2026. \url{https://oeis.org/A005258}.

\bibitem{cohen}
Henri Cohen, \emph{Ap\'ery Acceleration of Continued Fractions},
arXiv:2401.17720, 2024. In particular, Sections 3--4 discuss
staircase/diagonal acceleration and the examples $\zeta(2)$ and
$\zeta(3)$. \url{https://arxiv.org/abs/2401.17720}.

\bibitem{straub}
Armin Straub, \emph{Multivariate Ap\'ery numbers and supercongruences
of rational functions}, Algebra \& Number Theory \textbf{8} (2014),
no.~8, 1985--2008.
\href{https://doi.org/10.2140/ant.2014.8.1985}{doi:10.2140/ant.2014.8.1985}.
Preprint: \url{https://arxiv.org/abs/1401.0854}.

\bibitem{bacher}
Roland Bacher, Pierre de la Harpe, and Boris Venkov,
\emph{S\'eries de croissance et polyn\^omes d'Ehrhart associ\'es aux
r\'eseaux de racines}, Annales de l'Institut Fourier \textbf{49}
(1999), no.~3, 727--762.
\href{https://doi.org/10.5802/aif.1689}{doi:10.5802/aif.1689}.
\url{https://www.numdam.org/item/AIF_1999__49_3_727_0/}.

\end{thebibliography}
\end{document}
